\section{Analysephase}
\label{sec:Analysephase}
Ziel der Analysephase ist es, Anforderungen zu definieren, die bis zum Projektabschluss in der bestehenden \Abkuerzung{Java
EE}-Anwendung umgesetzt werden müssen. Dazu wird in der Ist-Analyse das bisherige System hinsichtlich der Task-Übersicht
untersucht und beschrieben. Im Anschluss schildert das Soll-Konzept, wie die Zielsetzung\VerweisVF{sub:Motivation und
Zielsetzung} mit dem Entwicklungsstand zum Zeitpunkt der Betrachtung erreicht werden soll. Nach einer kurzen Prüfung weiterer
zukünftiger Anforderungen wird abschließend eine für das Projektumfeld hinreichende Wirtschaftlichkeitsbetrachtung dargelegt.

\subsection{Ist-Analyse} 
\label{sub:Ist-Analyse}
Ein Task wird vom \Abkuerzung{VMS} automatisch generiert und ermöglicht bei bestimmten Ereignissen die Interaktion mit den
Benutzern. Diese Task sind in vordefinierte Task-Typen unterteilt: Fehler, Information und Benutzerinteraktion (z.B. Bestätigung
einer Vertragsänderung oder eines neuen Benutzers).

In der Task-Übersicht findet der Endanwender alle offenen Tasks sowie die in Bearbeitung befindendlichen Task des eigenen
Benutzers und seiner Benutzergruppen. Ein Benutzer kann nur aktiv nachsehen wie viele Tasks oder Info-Meldungen existieren. Dazu
muss er die Task-Übersicht öffnen und nach einem spezifischen Task-Typ filtern. Erst dann wird die Anzahl der Task-Einträge für
diesen Task-Typ ersichtlich. Der Benutzer wird weder auf neue Tasks hingewiesen, noch augenblicklich über diese informiert.
\citep[vgl.][]{Antrag}

\subsection{Soll-Konzept} 
\label{sub:Soll-Konzept}
Um einen schnellen Zugriff auf die Informationen "`Anzahl der offenen Tasks pro Typ"' gewährleisten zu können, sollen diese in die
Navigationsleiste neben dem Menüpunkt "`Tasks"' integriert werden. Die Darstellung erfolgt getrennt nach Task-Typ. Als
Darstellungsmittel werden einfache Symbole genutzt, welche die genannte Anzahl beinhalten. Dabei erhält jeder Typ eine Form und
eine Farbe. Als Begrifflichkeit dafür wird im weiterem Verlauf "`Status-Anzeige"' genutzt. Dies kann sich auf die Anzeige eines
oder aller Task-Typs beziehen.\\
Dem Kunden sind die Task-Typen "`Fehler"' und "`Benutzerinteraktion"' besonders wichtig. Die Status-Anzeige der Tasks soll sich
daher auf diese beiden beschränken.\\
Der Fachanwender kann außerdem visuell explizit auf neue Tasks hingewiesen werden. Des Weiteren können die jeweiligen
Status-Anzeigen auf die Task-Übersicht verlinkt werden. Die Übersicht wird dann entsprechend des Task-Typs gefiltert angezeigt.
Dabei darf der Benutzer-spezifische Filter in der Weboberfläche visuell nicht verändert werden und muss weiterhin mit dem
Benutzer-bekannten "`Filter anwenden"'-Button nutzbar sein. Die \AbbildungV{app:Mockup} dient zur Veranschaulichung dieser
Anforderung.

Zur Optimierung der Netzwerk- und Systemlast, muss darauf geachtet werden, dass die Anzeige vom \Abkuerzung{VMS} nur dann
aktualisiert wird, wenn sich die angezeigte Information serverseitig geändert hat. Das herkömmliche Verfahren des sog.
"`polling"' ist ungeeignet, da jeder Benutzer in einem festgelegtem Intervall die Informationen abfragt und damit eine erhöhte
Last auf dem System erzeugt. Daher soll unter Anderem auf bekannte Verfahren wie "`Comet"' oder auch "`long-polling"' gesetzt
werden. Bei diesen wird eine Anfrage zunächst vom Server nicht beantwortet sondern suspendiert. Liegt serverseitig ein Event vor,
wird diese angehaltene Verbindung entsprechend beantwortet. Erst dann stellt der Browser erneut eine Anfrage. Er erhält also nur
eine Antwort, wenn der Server dazu Bedarf sieht \citep[vgl.][]{HTTPLP}.

\begin{figure}[t]
\centering
\includegraphicsKeepAspectRatio{Bilder/mockup_status-anzeige.png}{1}
\caption{Mockup Status-Anzeige - Einbettung und Verknüpfung zur Task-Übersicht}
\label{app:Mockup} 
\end{figure}

\subsection{Zukünftig mögliche Anforderungen} 
\label{sub:Zukünftig mögliche Anforderungen}
Da es sich bei der \Abkuerzung{TSO BSS} um ein Langzeitprojekt handelt, werden zukünftig zu erwartende Anforderungen im Konzept
berücksichtigt. 

Denkbar ist beispielsweise, dass die hier eingeführte Event-Technologie auch für andere Zwecke verwendet wird. Es könnte \zB
gefordert werden, dass eine bestimmte Benutzergruppe direkten Einfluss auf einen Hintergrundprozess nehmen kann. Ein Prozess
teilt dem \Abkuerzung{VMS} mit, dass eine Benutzerinteraktion erforderlich ist. Daraufhin löst \Abkuerzung{VMS} ein Event aus,
dass bei einem Benutzer ein Popup anzeigt, welches die Aktion nachfragt. \\
Des Weiteren könnte sich der Kunde dahingehend umentscheiden, dass nicht nur die in \Verweis{sub:Soll-Konzept} genannten
Task-Typen angezeigt werden sollen, sondern auch weitere.\\
Als Letztes wird davon ausgegangen, dass neben dem letzten Besuch der Offenen-Task-Übersicht eines Benutzers in Zukunft weitere
Benutzer-abhängige Daten im \Abkuerzung{VMS} gesichert werden sollen. Dies wurde jedoch erst im Laufe der Realisierungsphase
spezifiziert und im Rahmen des Puffers für Anforderungsänderungen, in \VerweisVL{subsub:UserStatusRepository}, beachtet.

\subsection{Wirtschaftlichkeitsbetrachtung} 
\label{sub:Wirtschaftlichkeitsbetrachtung}
Zur Einhaltung des ökonomischen Prinzips wird die Wirtschaftlichkeit des Projektes überprüft. Dazu werden im Folgenden zunächst
die zu ermittelten Projektkosten mit dem erwarteten Nutzen verglichen. Abgeschlossen wird diese Betrachtung mit einer
Make-or-Buy-Entscheidung, die darlegt, warum nur eine Eigenentwicklung möglich ist und zur Realisierung ein kostenfreies Framework
genutzt wird.

\subsubsection{Kosten-Analyse} 
\label{subsub:Kosten-Analyse}
Die Kosten für das Projekt berechnen sich aus den Entwicklungskosten sowie den Kosten für zukünftige Anpassungen und Erweiterungen
des daraus entstandenen Features.

Zur Berechnung der Entwicklungskosten wird die in \VerweisKL{sub:Zeitlicher Rahmen des Projektes} genannte Projektdauer von 70
Stunden genutzt und ein interner Verrechnungssatz von \eur{45} pro Stunde\footnote{Ausgedachter Wert; Der interne Satz darf nicht
bekannt gegeben werden.} verwendet.
Dieser setzt sich aus verschiedenen Kostenfaktoren zusammen. Dazu zählen unter anderem: Lohn- und Lohnnebenkosten, Energie,
IT-Infrastruktur, Verwaltung, Nutzung von Räumlichkeiten \etc Somit ergibt sich für die reinen Entwicklungskosten ein Wert in Höhe
von:
\[ 70h * 45{\frac{\eur{}}{h}} = \eur{3150} \]

\subsubsection{Nutzen-Analyse} 
\label{subsub:Nutzen-Analyse}
In der Nutzen-Analyse wird aufgezeigt, welcher Nutzen sich für die \Abkuerzung{BTC} aus der Investition in das Projekt "`\titel"'
 ergibt. Dabei muss zwischen monetärem und nicht-monetärem Nutzen unterschieden werden.

\paragraph{Monetärer Vorteil}
Ein monetärer Vorteil entsteht im Regelfall durch die Verkaufserlöse einer Software. In Fall von \Abkuerzung{VMS} handelt es sich
um eine Individualsoftware für einen Kunden, die kontinuierlich im Auftrag des Kunden weiterentwickelt wird. Der monetäre Vorteil
erschließt sich also durch das Anbieten und gewinnbringende Verkaufen dieser Erweiterung in der nächsten Beauftragung. Der Erlös
wird mit einem Verrechnungssatz von \eur{145} pro Stunde\footnote{Ausgedachter Wert; Der reelle Satz darf nicht
bekannt gegeben werden.} ermittelt. Daraus kann der zu erwartene Gewinn wie folgt berechnet werden:
\[ \bigg(70h * 145{\frac{\eur{}}{h}}\bigg) - \eur{3150} = \eur{7000} \]
	
\paragraph{Nicht-monetärer Vorteil}
Ein nicht-monetärer Aspekt liegt vor, wenn dieser gar nicht bis unzureichend monetär bewertbar ist. Durch das Anbieten der
Erweiterung signalisiert die \Abkuerzung{BTC} gegenüber dem Kunden das Verständnis seiner Bedürfnisse. Dies führt zu einer
Verbesserung der Kundenzufriedenheit und -beziehung, welches wiederum zu weiteren Aufträgen führen kann. Das
Task-Benachrichtigungssystem führt auch zu einer Verbesserung des operativen Tagesgeschäfts beim Kunden. So können die Mitarbeiter
schneller über bestimmte Ereignisse informiert werden, was eine Optimierung der Reaktionszeiten bedeutet. Sobald in Oktober 2015
in der Energie-Branche stündliche Vermarktungen eingeführt werden, wird beim Kunden dieser Faktor an Gewicht gewinnen.

\subsubsection{Make-or-Buy-Entscheidung} 
\label{subsub:Make-or-Buy-Entscheidung}
Die in \VerweisVL{sub:Soll-Konzept} beschriebene Erweiterung ist so individuell, dass am Markt keine Lösung angeboten wird.
Selbst im Falle eines standarisierten Marktproduktes könnte dieses nicht ohne größere Anpassungen und Implementierungsarbeiten
eingesetzt werden. Da es sich außerdem bei \Abkuerzung{VMS} um eine ganzheitliche Eigenentwicklung der \Abkuerzung{BTC} handelt,
soll auch weiterhin auf eine betriebsinterne Entwicklung gesetzt werden.

\paragraph{Framework} Die Umsetzung der Kommunikation zwischen Server und Client sowie der Verbindungshaltung stellt einen
kritischen Faktor dar. Das Atmosphere Framework bietet dafür geeignete Verfahren gemäß \VerweisVL{sub:Soll-Konzept} an. Auch um im
Rahmen des magischen Dreiecks\footnote{Vgl.: \url{http://www.projektmanagementhandbuch.de/projektinitiierung/zieldefinition/}
(Abruf: 24.03.2015).}: "`Zeit, Qualität und Kosten"' zu bleiben, wird auf das frei erhältliche Framework gesetzt. Dieses reduziert
die benötigte Entwicklungszeit so, dass die geplante Projektdauer ohne Einschnitte in der Produktqualität eingehalten werden kann.
